\begin{textsample}{POS Dim 3 – pb – Score 28.00 – pb/pb\_inf\_6.txt}  \label{ex:f3_pos_083}
3.4.1 Conviver O direito de conviver aponta para a criação de vínculos, que devem ser pensados à luz da democracia e das diversas relações pelas quais as \textbf{crianças} também vão se constituindo como sujeitos : relações com os \textbf{pares}, com \textbf{adultos}, com toda a comunidade, escolar e mais ampla.
É importante frisar que a convivência é um processo que coloca em jogo as experiências e saberes que as \textbf{crianças} trazem consigo.
Seus pertencimentos, suas origens, suas diversas linguagens, seus jeitos próprios de estabelecer \textbf{laços}, de \textbf{brincar}, considerando como cada \textbf{criança} é, suas experiências anteriores, mas também as oportunidades para que ela amplie seu repertório, na direção de conviver respeitando as diferenças, a natureza, as diversas culturas e as diferenças entre as pessoas.
Ainda em relação à convivência, é necessário também frisar que o acolhimento da \textbf{criança} e da sua família é fundamental.
A convivência diz respeito às famílias também, que precisam estar próximas às propostas que os/as professores/as e \textbf{instituição}, de um modo geral, pensam para as \textbf{crianças}, dialogando com elas e colocando seus olhares e \textbf{sugestões}. 3.4.2 \textbf{Brincar} O \textbf{brincar} é um dos direitos que também devem ser pensados na relação com a democracia.
É também \textbf{brincando} que a \textbf{criança} estabelece relações com os outros e com o mundo, podendo, por meio delas, vivenciar situações que exigem respeito, reconhecimento ao pensamento do outro, dentre outras possibilidades.
Não há um único modo de \textbf{brincar} ou uma \textbf{brincadeira} melhor.
As \textbf{brincadeiras} que fazem parte do acervo cultural de um determinado grupo devem ser consideradas e reconhecidas.
As \textbf{brincadeiras} fazem parte de um patrimônio cultural, por um lado, mas também nascem das interações entre as \textbf{crianças} e entre elas e os \textbf{adultos}, nas condições e materialidades próprias ao tempo que lhes constituem, podendo ser invenções ou criações que instigam a imaginação, ampliam as culturas \textbf{infantis}, traduzem emoções e propiciam vivências \textbf{sensoriais}, corporais, cognitivas, sociais, subjetivas.
O \textbf{brincar} é uma experiência humana desde a antiguidade, uma vez que há evidências de que o homem sempre \textbf{brincou}.
Porém, durante muito tempo, essa atividade foi – e, de certa forma, ainda é – considerada sem importância.
Para muitos \textbf{adultos} – famílias e mesmo professores/as e \textbf{instituições} educacionais – o \textbf{brincar} e a \textbf{brincadeira} ainda são considerados algo secundário.
Estudos de diversos autores comprovam que as \textbf{crianças} aprendem melhor por meio das \textbf{brincadeiras}.
Não se trata apenas de aprendizagem no sentido cognitivo do termo, mas aprendizagem de relações, de respeito, da cultura do outro, dentre outras.
Sendo assim, o \textbf{brincar} deve ser reconhecido e vivenciado nas práticas das \textbf{instituições} de Educação \textbf{Infantil}, sem restrições à faixa \textbf{etária} das \textbf{crianças}.
Mas, como os/as professores/as da Educação \textbf{Infantil}, veem o \textbf{brincar}?
Como um direito?
É importante pensar sobre o lugar do \textbf{brincar}, da \textbf{brincadeira} e do \textbf{brinquedo} nas práticas desenvolvidas com as \textbf{crianças}.
E elas podem ser muitas e diversificadas, como a vivência com as palavras, as cantigas, os contos, os corpos, os jogos, com materiais diversos, elementos da natureza, com situações diversificadas.
Enfim, o repertório é infinito de possibilidades desde que o olhar sobre o \textbf{brincar} considere a potencialidade dessa prática na própria constituição da identidade da \textbf{criança}.
É importante, nessa direção, que o/a > A interação durante o \textbf{brincar} caracteriza o cotidiano da \textbf{infância}, trazendo \textbf{consigo} muitas aprendizagens e potenciais para o desenvolvimento integral das \textbf{crianças}. ( BNCC, 2017, EDUCAÇÃO \textbf{INFANTIL} ) professor/a pense em situações que contemplem a \textbf{brincadeira} relacionada à inclusão de \textbf{crianças} que têm pertencimentos culturais específicos e/ou deficiências, considerando -as nas práticas realizadas na \textbf{instituição} de Educação \textbf{Infantil}. 3.4.3 Participar A participação da \textbf{criança} na gestão da \textbf{instituição} educacional e das atividades cotidianas nela realizadas possibilita a esse sujeito desenvolver autonomia e tornar -se parte interessada nos destinos da \textbf{instituição} a qual pertence.
Isso requer participação ativa em todas as ações desenvolvidas na \textbf{instituição}, fazendo com que a \textbf{criança}, desde o \textbf{bebê}, não só defina destinos na \textbf{instituição}, mas se responsabilize por eles.
Escolher as \textbf{brincadeiras} que deseja participar, os materiais e espaços onde quer \textbf{brincar} ou realizar alguma ação favorece a implicação da \textbf{criança} com a \textbf{instituição} e um maior conhecimento sobre ela e sobre si mesma.
As eleições para escolha da direção das \textbf{instituições} podem ser um excelente momento para a \textbf{criança} exercer o direito de participar, mas também escolher coletivamente onde uma faixa deve ser exposta, que lanche pode escolher para comemorar alguma data especial, dentre outras possibilidades, contribuem para que a \textbf{criança} exerça esse direito e conheça mais sobre o funcionamento da \textbf{instituição} na qual está inserida.
A escolha das \textbf{brincadeiras}, o \textbf{manuseio} dos materiais utilizados, a escolha dos ambientes, tudo isso deve ser compartilhado com as \textbf{crianças} de modo que as mesmas possam opinar e participar da tomada de decisões.
Dessa forma, participando de experiências diversas, a \textbf{criança} aprende a lidar com conflitos, a compreender direitos e limitações e pode avançar positivamente em seus processos de aprendizagem e desenvolvimento como cidadã que é e deve ser tomada no espaço institucional. 3.4.4 Explorar A \textbf{instituição} de Educação \textbf{Infantil}, considerando a sua oferta em \textbf{creche} e \textbf{pré-escola}, deve estar \textbf{atenta} à relação da \textbf{criança} com o mundo pela via da exploração.
Se a \textbf{criança} participa ativamente da vida de sua \textbf{instituição}, ela estará disponível para descobrir, conhecer e criar pela ótica da exploração.
Assim, corpo, movimentos, \textbf{gestos}, \textbf{sons}, espaços diversos, situações, mudanças nos tempos e na organização da \textbf{rotina}, novidades, histórias \textbf{contadas} oral e livremente ou lidas a partir de um texto, objetos variados quanto ao tamanho, a cor, forma, espessura, enfim, constituem possibilidades infinitas de situações que podem e devem ser exploradas pelas \textbf{crianças}, com vistas a lhes favorecer o acesso ao mundo, conhecendo -o, mas também a interação entre as suas descobertas e o outro.
Com esse direito garantido, a \textbf{criança} potencializa a ampliação dos seus saberes, sua linguagem e seus conhecimentos.
Uma ação que pode fornecer caminhos para a exploração diz respeito ao próprio corpo da \textbf{criança}, que pode ser vivenciado pela \textbf{criança} por meio de músicas, de experiências com argila, com cores, na direção de que ela se autoconheça e perceba a si mesma e ao outro como sujeitos em interação.
Outra ação é explorar o ambiente onde a \textbf{criança} está inserida, tanto em sua casa como na \textbf{instituição} de Educação \textbf{Infantil}.
Nesse sentido, é importante saber sobre o pertencimento cultural da \textbf{criança}.
Ela reside no campo ou na cidade?
Mora perto de rios, cachoeiras, do mar?
Vive em apartamento ou casa?
Tem contato com a natureza?
O que as \textbf{crianças} podem aprender com a natureza?
O que aprendem neste contato direto e sensível?
O que pode ser realizado junto aos \textbf{bebês}, às \textbf{crianças} bem \textbf{pequenas} e \textbf{pequenas} para aproximar, por exemplo, a \textbf{criança} da natureza?
A relação com a natureza tem muito a contribuir na construção de saberes pela \textbf{criança}.
Seu contato com ela aproxima e constrói uma relação, na qual ela pode perceber possibilidades de convivência sensível, na direção de que a compreenda não como estranha, mas como um espaço do qual a \textbf{criança} faz parte, construindo e ampliando sua consciência pessoal, planetária e ecológica.
Podemos perceber apenas por esse exemplo que exercitar o direito de explorar contribui de forma significativa na construção e no desenvolvimento das \textbf{crianças} na sua relação com os outros e com o mundo. 3.4.5 Expressar A \textbf{criança} tem direito de expressar e expressar -se na \textbf{instituição} de Educação \textbf{Infantil} e em todos os demais espaços nos quais se situa.
E ela o faz por meio de diferentes linguagens que devem ser compreendidas como legítimas à sua comunicação : o choro, o silêncio, o \textbf{gesto}, a palavra, o grito, o sono, a agressividade, a alegria, a emoção, enfim, todas as formas de sentir e viver sua existência devem ser \textbf{acolhidas} como modos de expressão.
Nesse sentido, as diversas linguagens devem ser concebidas como meios de expressão da \textbf{criança} no espaço da \textbf{instituição} de Educação \textbf{Infantil} e ela deve ter direito a expor o que sente, o que pensa, o que cria/imagina, o que entende e não entende, o que gosta e não gosta.
Enfim, ela tem direito de comunicar aos outros, por meio de suas expressões, sentimentos, pensamentos, ideias, sensações. 3.4.6 Conhecer -se > CONHECER -SE e construir sua identidade pessoal, social e cultural, constituindo uma imagem positiva de si e de seus grupos de pertencimento, nas diversas experiências de \textbf{cuidados}, interações e \textbf{brincadeiras} vivenciadas na \textbf{instituição} de Educação \textbf{Infantil}.
A \textbf{criança} tem direito a conhecer -se e conhecer, potencializando a construção de uma imagem positiva de si e do seu grupo cultural e social.
A garantia deste direito contribui para a construção de sua identidade pessoal e coletiva, mas também, das diversas identidades das outras \textbf{crianças} e dos \textbf{adultos} com os quais ela se relaciona.
É na construção desse direito que a \textbf{criança} começa a perceber aspectos que a distingue das outras pessoas, cabendo ao/à professor/a e à \textbf{instituição} valorizar seus pertencimentos, suas características étnicas, culturais e pessoais.
No \textbf{dia} a \textbf{dia}, as \textbf{crianças} vão se percebendo e percebendo os outros como diferentes.
Essa \textbf{percepção} deve ser pensada e elaborada à luz do respeito, da democracia e da valorização das diferenças, favorecendo a construção de um conjunto de valores, crenças e conhecimentos diversificados, em que a prioridade seja valorizar e respeitar o outro como ele é. > Conhecer -se e construir sua identidade pessoal, social e cultural, constituindo uma imagem positiva de si e de seus grupos de pertencimento, nas diversas experiências de \textbf{cuidados}, interações, \textbf{brincadeiras} e linguagens vivenciadas na \textbf{instituição} escolar e em seu contexto familiar e comunitário. ( BNCC, 2017, EDUCAÇÃO \textbf{INFANTIL} ) > > Expressar, como sujeito dialógico, criativo e sensível, suas necessidades, emoções, sentimentos, dúvidas, hipóteses, descobertas, opiniões, questionamentos, por meio de diferentes linguagens. ( BNCC, 2017, EDUCAÇÃO \textbf{INFANTIL} ) Cabe aos/às professores/as \textbf{acolher} e possibilitar espaços para que as \textbf{crianças} possam expressar e expressar -se de diversos modos, manifestando a sua compreensão da importância da expressão dos seus sentimentos e descobertas como forma de dizer o que sente, esclarecer suas dúvidas sem medo, levantar hipóteses, confirmar ou não essas hipóteses, sem que sofram repressões.

% matched lemmas: acolher, adulto, atentar, bebê, brincadeira, brincar, brinquedo, conseguir, contar, creche, criança, cuidado, dia, etário, gesto, infantil, infância, instituição, laço, manuseio, par, pequeno, percepção, pré-escola, rotina, sensorial, som, sugestão
\end{textsample}
